% Clase 11 — planteo de la esfera dielectrica en campo uniforme.
% Lo que la figura tiene que mostrar: las DOS zonas (1 afuera, 2 adentro), el
% escalon de permitividad en r = a, y las coordenadas (r, theta) medidas desde
% el eje z, que es la direccion de E0. Las condiciones de borde van en la tabla
% del texto, no aca.
\begin{tikzpicture}[scale=1]

  % campo aplicado (lejos, a la izquierda)
  \foreach \yy in {-1.6,-0.8,0,0.8,1.6}{
    \draw[figvec] (-4.6,\yy) -- (-3.5,\yy);}
  \node[figlbl, figblue] at (-4.05,2.15) {$\vec E_0$};

  % esfera
  \fill[figamber!12] (0,0) circle (1.6);
  \draw[figamber, line width=0.9pt] (0,0) circle (1.6);

  % eje z
  \draw[figaxis] (-2.6,0) -- (3.4,0) node[right, figlbl, figgray] {$z$};

  % radio a
  \draw[figvecaux] (0,0) -- (250:1.6);
  \node[figlblaux, fill=figamber!12, inner sep=1pt] at (236:0.95) {$a$};

  % punto de campo y coordenadas
  \coordinate (O) at (0,0);
  \coordinate (P) at (40:2.8);
  \draw[figvec] (O) -- (P);
  \fill[figblue] (P) circle (1.6pt);
  \node[figlbl, figblue, fill=white, inner sep=1pt] at ($(O)!0.62!(P)+(0.05,0.3)$) {$r$};
  \draw[figgray, line width=0.6pt] (0.75,0) arc (0:40:0.75);
  \node[figlblaux] at (20:1.05) {$\theta$};

  % zonas
  \node[figlbl, figamber!80!black] at (-0.55,0.75) {zona 2};
  \node[figlblaux, figamber!80!black] at (-0.55,0.35) {$\varepsilon$};
  \node[figlbl] at (2.55,-1.55) {zona 1};
  \node[figlblaux] at (2.55,-1.95) {$\varepsilon_0$};

\end{tikzpicture}
